\documentclass[10pt,a4paper]{amsart}
\usepackage[margin=19mm]{geometry}
\usepackage{amsmath,amssymb,mathtools}
\usepackage[scr=rsfs,cal=euler]{mathalfa}
\usepackage{xfrac}
\usepackage[unicode,hidelinks,bookmarks=false]{hyperref}
\newcommand{\ie}{i.e.\ }
\newcommand{\cf}{cf.\ }
\newcommand{\resp}{respectively\ }
\newcommand{\Subsection}{\subsection}
\providecommand{\nobreakdash}{\nobreak\hbox{-}}
\setlength{\emergencystretch}{2em}
\multlinegap0pt
\usepackage{amssymb}
\usepackage{xcolor}
\usepackage[all]{xy}
\usepackage[shortlabels]{enumitem}
\newtheorem{lemma}{Lemma}
\DeclareMathOperator{\Hom}{Hom}
\newcommand{\sm}{\sigma[M]}
\title{Reduced rank in $\sm$}
\author{John A. Beachy, Mauricio Medina-B\'arcenas}
\date{Source: arXiv:2201.07196v3 (CC BY 4.0)}
\begin{document}
\maketitle

\noindent\emph{Source and licence.} Reduced rank in $\sm$. Version arXiv:2201.07196v3 (CC BY 4.0), distributed by arXiv under the Creative Commons Attribution 4.0 International licence, http://creativecommons.org/licenses/by/4.0/, as recorded in the arXiv licence metadata for this exact version and shown on the abstract page https://arxiv.org/abs/2201.07196v3. Full licence notice: \texttt{licenses/arxiv-2201.07196-CC-BY-4.0.txt}.\par\smallskip

\noindent\textit{Editorial scope.} Counted unit: the article's lemma lemma (source0.tex lines 367--377). The article's complete original proof of it is delivered with it (source0.tex lines 379--417, one \texttt{proof} environment of 39 lines). No mathematics is written, repaired or re-derived: the statement and the proof are the article's own lines, in the article's own notation, and no retained line is substituted or re-flowed.  Retained uncounted as the source-local context the proof needs: lines 359--365.  Nothing was omitted from any retained span, so the package carries no omission record.  No retained line contains a citation, so the wrapper carries no bibliography and no bibliography entry.  No retained line contains a figure environment, an image-inclusion command or drawing code, so no image is delivered and the package declares no asset dependency.  Editorial formatting only, in addition: an AMS \texttt{amsart} preamble that restates the article's own declarations and macros used by the retained lines, namely the environments \texttt{lemma} on separate counters, together with the standard packages those lines need.  The retained environments print in this document as Lemma 1 in order of appearance, whereas the article prints the same items with the numbering of its own full document.  The complete native source of the article as distributed in the arXiv e-print for this exact version is bundled unchanged as \texttt{sources/120965/source0.tex}.
	\begin{equation}\label{11}
		\begin{split}
			\mathcal{T}& =\{N\in\sm\mid \Hom_R(N,E)=0\} \\
			\mathcal{F}& =\{L\in\sm\mid L\hookrightarrow \textstyle{\prod_I^{[M]}E}\;
			\text{for some index set}\;I\}
		\end{split}
	\end{equation}

	\begin{lemma}
		Let $M$ be an $R$-module. 
		If $\tau=(\mathcal{T},\mathcal{F})$ 
		is a pair of classes of modules in $\sm$, 
		then $\tau$ is a hereditary torsion theory in $\sm$ 
		if and only if 
		there exists a hereditary torsion theory 
		$(\mathfrak{T},\mathfrak{F})$ in $R$-Mod 
		such that $\mathcal{T}=\mathfrak{T}\cap\sm$ and 
		$\mathcal{F}=\mathfrak{F}\cap\sm$.
	\end{lemma}

	\begin{proof}
		Suppose that $\tau=(\mathcal{T},\mathcal{F})$ 
		is a hereditary torsion theory in $\sm$. 
		Let $E$ be an $M$-injective module $E$ satisfying 
		the conditions of (\ref{11}), 
		and let $Q$ be the injective hull of $E$ in $R$-Mod. 
		Then $Q$ cogenerates a hereditary torsion theory 
		$(\mathfrak{T},\mathfrak{F})$ in $R$-Mod. 
		We claim that $\mathcal{T}=\mathfrak{T}\cap\sm$ and 
		$\mathcal{F}=\mathfrak{F}\cap\sm$. 
		Since $E\leq Q$, 
		it is clear that $\mathfrak{T}\cap\sm\subseteq \mathcal{T}$. 
		Suppose that there is a homomorphism $f:N\to Q$,
		for $N\in\mathcal{T}$. 
		Since $N$ is in $\sm$, 
		the image $f(M)$ is in $\sm$, 
		which implies that $f(M)\subseteq E$, and so $f=0$. 
		Thus $\mathcal{T}\subseteq\mathfrak{T}\cap\sm$. 
		On the other hand, since $\prod_I^{[M]}E\leq Q^I$ for any set $I$, 
		%and the product in $\sm$ is a submodule of the direct product in $R$-Mod
		we have that $\mathcal{F}\subseteq \mathfrak{F}\cap\sm$. 
		Now suppose that $N$ is in $\mathfrak{F}\cap\sm$, 
		so that there is a monomorphism $\alpha:N\to Q^I$. 
		For every index $i\in I$, 
		consider $\pi_i\alpha:N\to Q$, 
		where $\pi_i:Q^I\to Q$ is the canonical projection. 
		Since $N$ is in $\sm$ and the largest submodule of $Q$ in $\sm$ is $E$, 
		we have $\pi_i\alpha(N)\subseteq E$ for all $i\in I$. 
		This implies that $\alpha(N)\subseteq E^I$. 
		Since $\alpha(N)$ is in $\sm$, 
		$\alpha(N)$ must be contained in $\prod_I^{[M]}E$, 
		and thus $N$ is in $\mathcal{F}$. 
		
		Conversely, let $(\mathfrak{T},\mathfrak{F})$ 
		be a hereditary torsion theory in $R$-Mod. 
		Since $\sm$ is a hereditary pretorsion class,
		it follows easily that $(\mathfrak{T}\cap\sm,\mathfrak{F}\cap\sm)$ 
		is a hereditary torsion theory in $\sm$.	
	\end{proof}

\end{document}
